% Part: incompleteness
% Chapter: arithmetization-syntax
% Section: proofs-in-ax

\documentclass[../../../include/open-logic-section]{subfiles}

\begin{document}

\olfileid{inc}{art}{pax}
\olsection{مسلّمی \usetoken{P}{derivation}}

\begin{explain}
مسلّمی !!{derivation}s نوں عددی روپ دین لئی
سانوں !!{derivation}s نوں عدداں دے روپ وچ
ظاہر کرنا پینا اے۔ کیوں جے !!{derivation}s
بس !!{formula}s دے سلسلے نیں، ایس لئی
واضح طریقہ ایہ اے کہ ہر !!{derivation}
نوں اوس وچلے !!{formula}s دے کوڈاں دے
سلسلے دا کوڈ دیئیے۔
\end{explain}

\begin{defn}
جے $\delta$ اک مسلّمی !!{derivation} اے،
جیہڑا !!{formula}s $!A_1$، \dots،
~$!A_n$ توں بنیا اے، تاں
$\Gn{\delta}$ ایہ اے:
\[
\tuple{\Gn{!A_1}, \dots, \Gn{!A_n}}.
\]
\end{defn}

\begin{ex}
  ایہ سوکھا جیہا !!{derivation} ویکھو:
  \begin{derivation}
  1. & $!B \lif (!B \lor !A)$ \\
  2. & $(!B \lif (!B \lor !A)) \lif (!A  \lif (!B \lif (!B \lor !A)))$\\
  3. & $!A  \lif (!B \lif (!B \lor !A))$
  \end{derivation}
  ایس !!{derivation} دا گوڈل نمبر
  ایہ ہووے گا:
  \begin{align*}
  \openTuple\,
    & \Gn{!B \lif (!B \lor !A)}, \\
    &\Gn{(!B \lif (!B \lor !A)) \lif (!A  \lif (!B \lif (!B \lor !A)))},\\
    & \Gn{!A  \lif (!B \lif (!B \lor !A))} \,\closeTuple.
  \end{align*}
\end{ex}

\begin{explain}
!!{derivation}s دی عددی روپ طے کر کے سانوں ایہ وی دکھانا پینا اے
کہ اسیں انہاں !!{derivation}s اُتے
بنیادی بازگشتی ڈھنگ نال عمل کر
سکدے آں تے انہاں دیاں ضروری
خاصیتاں تے تعلقاں نوں وی ایس
طرح بیان کر سکدے آں۔ کجھ عمل
سوکھے نیں۔ مثال لئی، کسے
!!a{derivation} دے گوڈل نمبر~$d$
توں $(d)_{\len{d}-1}$ اوس دے
آخرلے !!{formula} دا گوڈل نمبر
دیندا اے۔ کجھ ہور عمل کافی
اوکھے نیں۔ اسیں گھٹ توں گھٹ
انہاں دا طریقہ خاکہ وار دساں
گے۔ مقصد ایہ دکھانا اے کہ
«$\delta$،~$!A$ دا
!!a{derivation} $\Gamma$ توں
اے» والا تعلق، $\delta$
تے~$!A$ دے گوڈل نمبراں دا بنیادی بازگشتی تعلق اے۔
\end{explain}

\begin{prop}
\ollabel{prop:followsby}
ایہ تعلق بنیادی بازگشتی نیں:
\begin{enumerate}
\item $!A$ اک مسلّمہ اے۔
\item $i$ویں سطر، جیہڑی $\delta$ وچ اے،
  دا جواز قاعدۂ اثباتِ مقدم نال اے۔
\item $i$ویں سطر، جیہڑی $\delta$ وچ اے،
  دا جواز \QR نال اے۔
\item $\delta$ اک درست
  !!{derivation} اے۔
\end{enumerate}
\end{prop}

\begin{proof}
سانوں دکھانا اے کہ !!{formula}s دے
گوڈل نمبراں تے !!{derivation}s دے
گوڈل نمبراں دے وچکار متعلقہ
تعلق بنیادی بازگشتی نیں۔
\begin{enumerate}
\item ساڈے کول مسلّمہ خاکیاں دی
  اک مقرر فہرست اے؛ تے $!A$
  مسلّمہ اوس ویلے اے جدوں اوہ
  کسے اک خاکے دے دتے روپ وچ
  ہووے۔ کیوں جے خاکیاں دی فہرست متناہی اے، ایس لئی
  ہر مسلّمہ خاکے لئی ایہ
  دکھانا کافی اے کہ اسیں
  بنیادی بازگشتی ڈھنگ نال
  پرکھ سکدے آں کہ $!A$
  اوس روپ وچ اے یا نہیں۔
  مثال لئی ایہ مسلّمہ
  خاکہ ویکھو:
  \[
  !B \lif (!C \lif !B).
  \]
  $!A$ ایس مسلّمہ خاکے دی
  مثال اوس ویلے اے جدوں
  !!{formula}s $!B$ تے $!C$
  ایہوجیہے ہون کہ $!A$
  ایویں ملے: '$($' نال
  $!B$، فیر '$\lif$'،
  فیر '$($'، فیر $!C$،
  فیر '$\lif$'، فیر $!B$
  تے آخر~'$))$' جوڑیا
  جائے۔ اسیں گوڈل نمبر $n$ دی،
  جیہڑا~$!A$ دا اے،
  متعلقہ خاصیت پرکھ سکدے
  آں، کیوں جے سلسلیاں
  نوں جوڑنا بنیادی
  بازگشتی اے، تے $!B$
  تے $!C$ دے گوڈل نمبر
  ضرور~$!A$ دے گوڈل
  نمبر توں چھوٹے نیں۔
  ایس دی وجہ ایہ اے کہ
  جدوں ایہ تعلق سچ
  ہووے، تاں $!B$
  تے $!C$ دونوں
  ~$!A$ دے ذیلی
  !!{formula}s نیں۔
  لہٰذا اسیں ایہ
  تعریف کر سکدے آں:
  \begin{multline*}
  \fn{IsAx}_{!B \lif (!C \lif !B)}(n) \defiff \bexists{b< n}{
    \bexists{c < n}{(\fn{Sent}(b) \land \fn{Sent}(c) \land {}}}\\ n =
  \Gn{(} \concat b \concat \Gn{\lif} \concat \Gn{(} \concat c \concat
  \Gn{\lif} \concat b \concat \Gn{))}).
  \end{multline*}
  جے ہر مسلّمہ خاکے
  لئی ایہوجیہی تعریف
  ہووے، تاں اوہناں دا منطقی «یا»
  خاصیت $\fn{IsAx}(n)$
  دی تعریف کردا اے:
  «$n$ کسے مسلّمے
  دا گوڈل نمبر اے»۔
\item $i$ویں سطر، جیہڑی $\delta$ وچ اے،
  دا جواز قاعدۂ
  اثباتِ مقدم نال
  عین اوس ویلے اے
  جدوں سطراں $j$
  تے $k < i$
  موجود ہون؛ سطر~$j$
  دا !!{sentence}
  کوئی فارمولا $!A$
  ہووے، سطر~$k$
  دا جملہ
  $!A \lif !B$
  ہووے، تے
  سطر~$i$ دا
  جملہ~$!B$
  ہووے۔
  \begin{multline*}
    \fn{MP}(d, i) \defiff \bexists{j < i}{\bexists{k < i}{}}\\
    (d)_k = \Gn{(} \concat (d)_j
      \concat \Gn{\lif} \concat (d)_i \concat \Gn{)}
  \end{multline*}
  کیوں جے محدود
  مسوّریت، سلسلیاں
  دا جوڑ تے $=$
  بنیادی بازگشتی
  نیں، ایہ اک
  بنیادی بازگشتی
  تعلق دی تعریف اے۔
\item $\delta$ دی کسے
  سطر دا جواز \QR{}
  نال اوس ویلے اے
  جدوں اوہ
  $!B \lif \lforall[x][!A(x)]$
  دے روپ وچ ہووے،
  پہلاں آ چکی کوئی سطر $!B \lif !A(c)$ ہووے، جتھے کوئی
  !!{constant}~$c$
  ہووے، تے $c$
  ~$!B$ وچ نہ
  آوندا ہووے۔
  ایہ عین اوس
  ویلے اے جدوں
  \begin{enumerate}
  \item کوئی !!a{sentence}~$!B$
    ہووے، تے
  \item !!a{formula}~$!A(x)$
    ہووے جیہدے وچ
    صرف متغیر~$x$
    آزاد ہووے؛ نال
    ایہ کہ
  \item سطر $i$ وچ
    $!B \lif \lforall[x][!A(x)]$
    ہووے،
  \item کوئی سطر
    $j < i$،
    $!B \lif \Subst{!A}{c}{x}$
    رکھدی ہووے،
    جتھے $c$ مستقل
    ہووے،
  \item تے اوہ
    ~$!B$ وچ نہ
    آوندا ہووے۔
  \end{enumerate}
  ایہ ساریاں گلّاں
  بنیادی بازگشتی
  ڈھنگ نال پرکھیاں
  جا سکدیاں نیں،
  کیوں جے $!B$،
  $!A(x)$ تے $x$
  دے گوڈل نمبر،
  سطر~$i$ دے
  فارمولے دے
  گوڈل نمبر توں
  چھوٹے نیں؛ تے
  $a$ دا گوڈل
  نمبر، سطر~$j$
  دے فارمولے دے
  گوڈل نمبر توں
  چھوٹا اے:
  \begin{multline*}
    \fn{QR}_1(d, i) \defiff \bexists{b < (d)_i}{\bexists{x <
        (d)_i}{\bexists{a < (d)_i}{\bexists{c < (d)_j}{(}}}}
    \\ \fn{Var}(x) \land \fn{Const}(c) \land {} \\
    (d)_i = \Gn{(} \concat
    b \concat \Gn{\lif} \concat \Gn{\lforall} \concat x \concat a
    \concat \Gn{)} \land {}\\
    (d)_j = \Gn{(} \concat b \concat
    \Gn{\lif} \concat \fn{Subst}(a,c,x) \concat \Gn{)} \land {}\\ \fn{Sent}(b)
    \land \fn{Sent}(\fn{Subst}(a,c,x)) \land {} \bforall{k <
      \len{b}}{(b)_k \neq (c)_0})
  \end{multline*}
  ایتھے اسیں $c$ نوں ثابتہ نشان دا تے $x$ نوں متغیر دا
  گوڈل نمبر سمجھدے آں، جدوں دوناں نوں ٹرم وجوں ویکھیے
  (یعنی انہاں دیاں علامتاں دے کوڈ نہیں)۔
  اسیں ایہ
  پرکھدے آں
  کہ $x$ ہی
  ~$!A(x)$ دا
  اکیلا آزاد
  متغیر اے؛
  ایس لئی
  $\Subst{!A(x)}{c}{x}$
  دا !!a{sentence}
  ہونا پرکھدے
  آں۔ تے ایہ
  یقینی بناندے
  آں کہ $c$،
  $!B$ وچ نہ
  آوے؛ ایس لئی
  منگدے آں
  کہ~$!B$ دی
  ہر علامت
  ~$c$ توں
  وکھری ہووے۔

  \QR{} دا
  دوجا روپ
  مشق وجوں چھڈدے آں۔
\item $d$ کسے
  درست !!{derivation}
  دا گوڈل نمبر
  عین اوس ویلے
  اے جدوں اوہدی
  ہر سطر یا
  تاں مسلّمہ
  ہووے، یا
  قاعدۂ اثباتِ
  مقدم نال، یا
  ~\QR نال
  جواز رکھدی
  ہووے۔ ایس لئی:
  \[
  \fn{Deriv}(d) \defiff \bforall{i < \len{d}}{(\fn{IsAx}((d)_i) \lor
    \fn{MP}(d,i) \lor \fn{QR}(d, i))}
  \]
\end{enumerate}
\end{proof}

\begin{prob}
ایہ تعلق
\olref[inc][art][pax]{prop:followsby}
دی طرز اُتے
تعریف کرو:
\begin{enumerate}
\item $\fn{IsAx}_{!A \lif (!B \lif (!A \land !B))}(n)$،
\item $\fn{IsAx}_{\lforall[x][!A(x)] \lif !A(t)}(n)$،
\item $\fn{QR}_{2}(d, i)$
  (\QR دے دوجے روپ لئی)۔
\end{enumerate}
\end{prob}

\begin{prop}
فرض کرو $\Gamma$،
!!{sentence}s دا
بنیادی بازگشتی
سیٹ اے۔ تاں
تعلق
$\Prf[\Gamma](x, y)$،
جیہڑا ایہ بیان
کردا اے کہ
«$x$ کسے
!!a{derivation}~$\delta$
دا کوڈ اے
جیہڑا $!A$ نوں~$\Gamma$ توں کڈھدا اے، تے
$y$،~$!A$ دا
گوڈل نمبر اے»،
بنیادی بازگشتی
اے۔
\end{prop}

\begin{proof}
فرض کرو
«$y \in \Gamma$»
بنیادی بازگشتی
محمول~$R_\Gamma(y)$ نال دِتّا گیا اے۔
سانوں دکھانا اے
کہ تعلق
$\Prf[\Gamma](x, y)$
بنیادی بازگشتی
اے؛ ایتھے
$\Prf[\Gamma](x,
y)$ عین اوس
ویلے سچ اے
جدوں $y$
کسے
!!a{sentence}~$!A$
دا گوڈل نمبر
ہووے تے $x$،
~$!A$ دے
!!a{derivation}
دا کوڈ ہووے
جیہڑا $\Gamma$
توں اے۔

پچھلے قضیے مطابق خاصیت
$\fn{Deriv}(x)$،
جیہڑی عین اوس
ویلے سچ اے
جدوں $x$
کسے درست
!!{derivation}~$\delta$
دا کوڈ ہووے،
بنیادی بازگشتی
اے۔ پر اوس
تعریف وچ سیٹ
$\Gamma$ نوں
!!{derivation}
دیاں سطراں
دا جواز دین
دا اک ہور
طریقہ نہیں
گِنّیا گیا۔
ساڈے بنیادی بازگشتی پرکھ وچ، جیہڑی ویکھدی اے
کہ سطر دا
جواز \QR{}
نال اے یا
نہیں، ایہ
شرط وی
نظرانداز سی
کہ مستقل~$c$
نوں~$\Gamma$
وچ آون دی
اجازت نہیں۔
ساڈی تعریف
نوں بدل کے
براہِ راست
$\Gamma$ دا
لحاظ رکھیا
جا سکدا اے،
پر $\fn{Deriv}$
تے استنباط
دا قضیہ
ورتنا سوکھا
اے۔
$\Gamma \Proves !A$
عین اوس
ویلے اے
جدوں
!!{sentence}s
دی کوئی متناہی فہرست $!B_1$،
\dots،
$!B_n \in \Gamma$
ایہوجیہی ہووے
کہ
$\{!B_1, \dots, !B_n\} \Proves !A$۔
تے استنباط
دے قضیے
مطابق ایہ
اوس ویلے اے
جدوں
$\Proves (!B_1 \lif (!B_2 \lif
\cdots (!B_n \lif !A)\cdots))$
ہووے۔
گوڈل نمبر~$z$
والا
!!a{sentence}
ایس روپ وچ
اے یا نہیں،
ایہ بنیادی
بازگشتی ڈھنگ
نال پرکھیا
جا سکدا اے۔
سو $x$ نوں
گوڈل نمبر~$y$
والے
!!{sentence}
دے
!!a{derivation}
دا گوڈل نمبر
\emph{$\Gamma$ توں}
سمجھّن دی
تھاں، اسیں
$x$ نوں
اُتّے دتے
روپ دے تہہ
در تہہ
شرطیہ دا
!!a{derivation}
دا گوڈل نمبر
سمجھدے آں،
جیہڑا~$\emptyset$
توں اے۔

پہلاں، جے
ساڈے کول
!!{sentence}s
دا سلسلہ
ہووے، تاں
اسیں بنیادی
بازگشتی ڈھنگ
نال ایہ
شرطیہ بنا
سکدے آں
جیہدے سارے
مقدم ایہ
جملے ہون
تے تالی
اک دتا
!!{sentence}
ہووے:
\begin{align*}
  \fn{hCond}(s, y, 0) & = y \\
  \fn{hCond}(s, y, n+1) & = \Gn{(} \concat (s)_{n} \concat \Gn{\lif}
  \concat \fn{Cond}(s, y, n) \concat \Gn{)}\\
  \fn{Cond}(s, y) & = \fn{hCond}(s, y, \len{s})\\
  \intertext{سو اسیں $\Prf[\Gamma](x, y)$ دی تعریف ایویں کر سکدے آں}
  \Prf[\Gamma](x, y) & \defiff \bexists{s < \fn{sequenceBound}(x,x)}{(} \\
  &\qquad (x)_{\len{x}-1} = \fn{Cond}(s,y) \land {} \\
  &\qquad\bforall{i<\len{s}}{(s)_i \in \Gamma} \land {} \\
  &\qquad\fn{Deriv}(x)).
\end{align*}
~$s$ دی حد
ایس گل توں ملدی اے کہ ہر
$(s)_i$،
!!{derivation}
دی آخری
سطر دے کسے
ذیلی
!!{formula}
دا گوڈل نمبر
اے؛ یعنی
اوہ
$(x)_{\len{x}-1}$
توں چھوٹا اے۔
مقدمیاں
$!B \in \Gamma$
دی گنتی،
یعنی~$s$
دی لمبائی،
~$x$ دی
آخری سطر
دی لمبائی
توں گھٹ اے۔
\end{proof}

\end{document}
